\chapter{Discussion}
\label{ch:discussion}

\section{Interpretation of Results}

\subsection{Where the Model Succeeds}

The hybrid framework demonstrates significant improvements over physics-only baselines in several areas:
\begin{itemize}
    \item \textbf{Accuracy}: Reduced error in energy and arrival time predictions, particularly for complex control strategies
    \item \textbf{Calibration}: Better-calibrated uncertainty estimates, enabling reliable risk assessment
    \item \textbf{Long-range dependencies}: Attention mechanisms successfully model force propagation through long consists
\end{itemize}

% TODO: Provide specific examples from results
% TODO: Discuss which scenarios show largest improvements

\subsection{Where the Model Fails}

Limitations observed in evaluation include:
\begin{itemize}
    \item \textbf{Extreme OOD scenarios}: Performance degrades on very long consists (200+ cars) or extreme grades (\SI{>3}{\percent})
    \item \textbf{Rare events}: Underestimation of tail probabilities for very high coupler forces
    \item \textbf{Computational cost}: While faster than high-fidelity simulation, transformer inference is still slower than pure physics models
\end{itemize}

% TODO: Quantify failure modes with examples
% TODO: Discuss potential mitigations

\section{Implications for Rail Labs and Planning Workflows}

\subsection{Decision Support Framing}

The framework enables rapid scenario exploration for:
\begin{itemize}
    \item \textbf{Energy planning}: Evaluating fuel consumption across different routes and control strategies
    \item \textbf{Safety analysis}: Identifying scenarios with high coupler force risk
    \item \textbf{Consist optimization}: Testing different car orderings and lengths
\end{itemize}

Probabilistic outputs support risk-aware decision making, allowing planners to choose strategies that balance efficiency and safety.

% TODO: Provide workflow examples
% TODO: Discuss integration with existing tools

\subsection{Uncertainty-Aware Planning}

The calibrated uncertainty estimates enable:
\begin{itemize}
    \item Setting appropriate safety margins based on predicted exceedance probabilities
    \item Identifying scenarios requiring high-fidelity simulation for validation
    \item Optimizing control plans while respecting probabilistic constraints
\end{itemize}

% TODO: Example of uncertainty-aware optimization
% TODO: Compare to deterministic planning approaches

\section{Limitations and Risk Management}

\subsection{Out-of-Distribution Generalization}

The model's performance degrades on scenarios far outside the training distribution. This limitation must be managed through:
\begin{itemize}
    \item \textbf{Domain awareness}: Identifying when inputs fall outside training distribution
    \item \textbf{Fallback strategies}: Using high-fidelity simulation for flagged scenarios
    \item \textbf{Active learning}: Expanding training data based on model uncertainty
\end{itemize}

% TODO: Discuss OOD detection methods
% TODO: Propose active learning framework

\subsection{Uncertainty Misuse}

There is a risk that users may misinterpret uncertainty estimates or use them inappropriately:
\begin{itemize}
    \item \textbf{Overconfidence}: Treating predictions as ground truth without considering uncertainty
    \item \textbf{Underconfidence}: Being overly conservative due to wide prediction intervals
    \item \textbf{Calibration drift}: Model calibration may degrade over time as operating conditions change
\end{itemize}

Mitigation requires clear documentation, user training, and periodic recalibration.

% TODO: Propose guidelines for uncertainty interpretation
% TODO: Discuss recalibration strategies

\section{Future Research Directions}

Several research directions could address current limitations:
\begin{itemize}
    \item \textbf{Real-world calibration}: Validating and improving calibration on operational data
    \item \textbf{Distributed power}: Extending to consists with multiple locomotives
    \item \textbf{Tighter safety constraints}: Incorporating regulatory requirements directly into the model
    \item \textbf{Integration with optimization}: End-to-end learning for control plan optimization
\end{itemize}

% TODO: Expand on each research direction
% TODO: Reference related work that could be built upon
